\documentclass{standalone}
\usepackage{circuitikz}
\usetikzlibrary{calc}
\definecolor{circuitnoteink}{HTML}{000000}
\definecolor{circuitnotemark}{HTML}{FE00FE}
\begin{document}
\begin{circuitikz}[american, line width=0.8pt]
\ctikzset{bipoles/length=1.2cm}
\coordinate (f3) at (2.4,-6);
\coordinate (f6) at (6,-6);
\coordinate (f9) at (9.6,-6);
\coordinate (f12) at (13.2,-6);
\coordinate (f15) at (16.8,-6);
\coordinate (f18) at (20.4,-6);
\node[dipchip, num pins=8, font=\scriptsize] (part-U1) at (f3) {}; % line 3
\node[dipchip, num pins=14, font=\scriptsize] (part-U2) at (f6) {}; % line 4
\node[dipchip, num pins=16, font=\scriptsize] (part-U3) at (f9) {}; % line 5
\node[dipchip, num pins=20, font=\scriptsize] (part-U4) at (f12) {}; % line 6
\node[dipchip, num pins=28, font=\scriptsize] (part-U5) at (f15) {}; % line 7
\node[dipchip, num pins=40, font=\scriptsize] (part-U6) at (f18) {}; % line 8
\draw[circuitnoteink] (0.6,-0.12) -- (22.2,-0.12); % line 10
\draw[circuitnoteink] (0.6,-12.12) -- (22.2,-12.12); % line 11
\draw[circuitnoteink] (0.6,-0.12) -- (0.6,-12.12); % line 12
\draw[circuitnoteink] (4.2,-0.12) -- (4.2,-12.12); % line 13
\draw[circuitnoteink] (7.8,-0.12) -- (7.8,-12.12); % line 14
\draw[circuitnoteink] (11.4,-0.12) -- (11.4,-12.12); % line 15
\draw[circuitnoteink] (15,-0.12) -- (15,-12.12); % line 16
\draw[circuitnoteink] (18.6,-0.12) -- (18.6,-12.12); % line 17
\draw[circuitnoteink] (22.2,-0.12) -- (22.2,-12.12); % line 18
\path (0.559,-1.073) rectangle (4.242,0.113); % line 19
\node[anchor=west, inner sep=0, circuitnotemark, font=\large] at (2.4,-0.48) {X}; % line 19
\path (1.003,-12.233) rectangle (3.797,-11.047); % line 20
\node[anchor=west, inner sep=0, circuitnotemark, font=\large] at (2.4,-11.64) {X}; % line 20
\path (4.032,-1.073) rectangle (7.969,0.113); % line 21
\node[anchor=west, inner sep=0, circuitnotemark, font=\large] at (6,-0.48) {X}; % line 21
\path (4.476,-12.233) rectangle (7.524,-11.047); % line 22
\node[anchor=west, inner sep=0, circuitnotemark, font=\large] at (6,-11.64) {X}; % line 22
\path (7.632,-1.073) rectangle (11.569,0.113); % line 23
\node[anchor=west, inner sep=0, circuitnotemark, font=\large] at (9.6,-0.48) {X}; % line 23
\path (8.076,-12.233) rectangle (11.124,-11.047); % line 24
\node[anchor=west, inner sep=0, circuitnotemark, font=\large] at (9.6,-11.64) {X}; % line 24
\path (11.232,-1.073) rectangle (15.168,0.113); % line 25
\node[anchor=west, inner sep=0, circuitnotemark, font=\large] at (13.2,-0.48) {X}; % line 25
\path (11.549,-12.233) rectangle (14.851,-11.047); % line 26
\node[anchor=west, inner sep=0, circuitnotemark, font=\large] at (13.2,-11.64) {X}; % line 26
\path (14.832,-1.073) rectangle (18.769,0.113); % line 27
\node[anchor=west, inner sep=0, circuitnotemark, font=\large] at (16.8,-0.48) {X}; % line 27
\path (15.149,-12.233) rectangle (18.451,-11.047); % line 28
\node[anchor=west, inner sep=0, circuitnotemark, font=\large] at (16.8,-11.64) {X}; % line 28
\path (18.432,-1.073) rectangle (22.368,0.113); % line 29
\node[anchor=west, inner sep=0, circuitnotemark, font=\large] at (20.4,-0.48) {X}; % line 29
\path (18.749,-12.233) rectangle (22.051,-11.047); % line 30
\node[anchor=west, inner sep=0, circuitnotemark, font=\large] at (20.4,-11.64) {X}; % line 30
\node[anchor=south west, inner sep=2pt, circuitnotemark, font=\LARGE\bfseries] (circuittitle) at (current bounding box.north west) {X};
\path (circuittitle.south west) rectangle ++(9.571,0.853);
\end{circuitikz}
\end{document}
